\documentclass[a4paper]{jsarticle}
\usepackage{gaiyo2}

%\documentclass[a4paper]{jarticle}
% として jarticle を使う場合は，上の \usepackage の部分を
%\usepackage[jslayout]{gaiyo2}
% のように [jslayout] を追加する．

% 以下は修論のフォーマットを参照して記入する．
\title{修論のタイトル}
\author{自分の名前}
\date{2024.3} % 例. 2024年3月卒業の場合は「2024.3」を記入する．
\kei{代数幾何解析数理科学系} % 自分の系を記入する．
\shisho{名前 教授} % 教授・准教授・助教の別に注意

\begin{document}
\begin{gaiyo}
  指導教員の指導のもとに修士論文の概要を作成する．概要の形式と
  内容は指導教員とよく相談するべきである．
  修士論文概要は，修士論文
  審査会に際して各系の教員が参照しながら発表を聞く．
  従って，各系の中である程度広い範囲で理解されるように書く
  必要がある．
  ここでは参考のために一般的な形式を紹介する．

  まず修士論文の基本的な内容を短く紹介する．これは
  系のメンバーが広く理解できるように書く．次に，修士論文の内容に関する背景を
  簡潔に紹介する．これは関連分野に属していれば理解
  できるように書く．続いて，修士論文の対象とした一般的な問題
  を述べる．
  修士論文に独自の結果を含むか否か，総合報告であるかどうかについて，
  ここまでに明示しておく．

  以上は主要な結果を述べるための準備である．以下，
  主要な結果を要約する．主定理を正確に記述し，必要な定義などを記述する．
  次に，修士論文に従って，主要な結果の証明の方針などを記述する．
  また，修士論文の章立て，各章の内容を簡潔に記述する．
  必要であれば，最後に参考文献を挙げる．
\end{gaiyo}
\end{document}
